% Four families of TTS systems placed on a common pipeline. Each box spans the
% stages that one trained model covers; arrows mark the interface between models.
\begin{tikzpicture}[
  x=1mm,y=1mm,figure picture,
  stage/.style={font=\figurefont\scriptsize\bfseries,text=FigureInk,align=center},
  fam/.style={font=\figurefont\scriptsize,text=FigureInk,anchor=east,align=right,text width=22mm},
  tm/.style={figure node,minimum height=9mm,inner ysep=2pt,font=\figurefont\scriptsize},
  front/.style={tm,draw=FigureInput,fill=FigureInputFill},
  acou/.style={tm,draw=FigureModel,fill=FigureModelFill},
  mid/.style={tm,draw=FigureStroke,fill=white},
  voc/.style={tm,draw=FigureLoss,fill=FigureLossFill},
  flow/.style={figure flow}
]
  % column headings
  \node[stage] at (14,8) {文本前端};
  \node[stage] at (55,8) {对齐与声学生成};
  \node[stage] at (93,8) {中间表示};
  \node[stage] at (124,8) {波形};
  \draw[figure link,draw=FigureGrid] (0,4) -- (140,4);

  % row 1: two-stage
  \node[fam] at (-2,-6) {两阶段流水线\\（\ref{sec:nlp-speech-acoustics}、\ref{sec:nlp-speech-waveform}）};
  \node[front,text width=24mm] (f1) at (14,-6) {规范化与G2P};
  \node[acou,text width=40mm] (a1) at (55,-6) {Tacotron 2：注意力隐式对齐\\FastSpeech 2：显式时长};
  \node[mid,text width=18mm] (i1) at (93,-6) {梅尔频谱};
  \node[voc,text width=26mm] (v1) at (124,-6) {WaveNet或\\HiFi-GAN};
  \draw[flow] (f1) -- (a1); \draw[flow] (a1) -- (i1); \draw[flow] (i1) -- (v1);

  % row 2: VITS
  \node[fam] at (-2,-20) {端到端VAE\\（\ref{sec:nlp-speech-vits}）};
  \node[front,text width=24mm] (f2) at (14,-20) {音素};
  \node[acou,text width=106mm] (a2) at (85,-20) {VITS：单调对齐搜索与随机时长；规范化流先验；隐变量经对抗训练的生成器直接得到波形};
  \draw[flow] (f2) -- (a2);

  % row 3: diffusion / flow matching
  \node[fam] at (-2,-34) {扩散与流匹配\\（\ref{sec:nlp-speech-diffusion-tts}）};
  \node[front,text width=24mm] (f3) at (14,-34) {音素或字符};
  \node[acou,text width=40mm] (a3) at (55,-34) {Grad-TTS：分数去噪\\F5-TTS：流匹配填充};
  \node[mid,text width=18mm] (i3) at (93,-34) {梅尔频谱};
  \node[voc,text width=26mm] (v3) at (124,-34) {声码器};
  \draw[flow] (f3) -- (a3); \draw[flow] (a3) -- (i3); \draw[flow] (i3) -- (v3);

  % row 4: codec language model
  \node[fam] at (-2,-48) {编解码词元\\语言模型（\ref{sec:nlp-speech-codec-lm-tts}）};
  \node[front,text width=24mm] (f4) at (14,-48) {音素$+$3秒提示};
  \node[acou,text width=40mm] (a4) at (55,-48) {VALL-E：首层自回归\\其余层非自回归};
  \node[mid,text width=18mm] (i4) at (93,-48) {离散声学\\词元};
  \node[voc,text width=26mm] (v4) at (124,-48) {编解码器\\解码器};
  \draw[flow] (f4) -- (a4); \draw[flow] (a4) -- (i4); \draw[flow] (i4) -- (v4);
\end{tikzpicture}
